\section{Ход работы}

\subsection{Наивный поиск (пункт~\ref{task:naive})}
Наивный алгоритм прикладывает образец к каждой позиции текста и сравнивает символы слева
направо до первого несовпадения. Реализация приведена в листинге~\ref{lst:naive}.

\bbtCodeFile[python, lines=21-34]{materials/code/lab05.py}{Наивный поиск со счётчиком
сравнений}{lst:naive}

Внешний цикл перебирает позиции от 0 до $n - m$, внутренний сравнивает символы текста и
образца и прерывается на первом несовпадении. Каждое сравнение увеличивает счётчик
\texttt{comparisons}. В худшем случае на каждой позиции совпадают все символы, кроме
последнего, и число сравнений равно
\begin{equation}\label{eq:naive}
  C = (n - m + 1)\, m,
\end{equation}
где $n$ — длина текста;\\
$m$ — длина образца.

Если длина образца пропорциональна длине текста, как в этой работе, оценка
\eqref{eq:naive} растёт как $n^2$.

\subsection{Алгоритм Кнута — Морриса — Пратта (пункт~\ref{task:kmp})}
Алгоритм КМП заранее вычисляет префикс-функцию образца (листинг~\ref{lst:prefix}) и при
несовпадении сдвигает образец сразу на нужное число позиций, не возвращаясь назад по тексту.

\bbtCodeFile[python, lines=37-48]{materials/code/lab05.py}{Вычисление префикс-функции}{lst:prefix}

Значение \texttt{pi[i]} — длина наибольшего собственного префикса строки
\texttt{pattern[:i + 1]}, который одновременно является её суффиксом. Переменная \texttt{k}
хранит длину текущего совпадения и при несовпадении уменьшается по уже вычисленным
значениям \texttt{pi}, поэтому префикс-функция вычисляется за $O(m)$.

Сам поиск приведён в листинге~\ref{lst:kmp}.

\bbtCodeFile[python, lines=51-70]{materials/code/lab05.py}{Поиск по алгоритму КМП}{lst:kmp}

Индекс \texttt{i} по тексту только растёт. При несовпадении меняется лишь \texttt{k} — длина
совпавшего начала образца, поэтому на каждый символ текста приходится не больше двух
сравнений в среднем, а весь поиск занимает $O(n + m)$. Найденное вхождение записывается, и
поиск продолжается с \texttt{k = pi[k - 1]}, поэтому перекрывающиеся вхождения не теряются.

\subsection{Входные данные и замеры (пункт~\ref{task:bench})}
Входы строит функция \texttt{make\_inputs} (листинг~\ref{lst:inputs}): случайный текст над
алфавитом \texttt{abcd} с образцом, вырезанным из этого же текста, и худший для наивного
поиска случай — текст из одних букв \texttt{a} и образец из букв \texttt{a} с буквой
\texttt{b} в конце. Длина образца — $n / 20$.

\bbtCodeFile[python, lines=73-85]{materials/code/lab05.py}{Генерация классов входов}{lst:inputs}

Время измеряет функция \texttt{median\_ms} (листинг~\ref{lst:median}): один прогревочный
запуск, затем пять замеров таймером \texttt{time.perf\_counter} и медиана.

\bbtCodeFile[python, lines=88-96]{materials/code/lab05.py}{Медиана времени поиска}{lst:median}

Замеры выполнены для длин текста 1000, 2000, 4000, 8000 и 16~000 символов по seed 107
(вариант 7) на Python 3.12.16. Медианы времени приведены в таблице~\ref{tab:time}, число
сравнений — в таблице~\ref{tab:comparisons}.

\begin{bbttable}{Медиана времени поиска, мс}{tab:time}
  \begin{tabular}{|r|c|c|c|c|}
    \hline
    & \multicolumn{2}{c|}{Случайный текст} & \multicolumn{2}{c|}{Худший случай} \\ \cline{2-5}
    $n$ & наивный & КМП & наивный & КМП \\ \hline
    1000    & 0,126 & 0,081 & 2,798   & 0,177 \\ \hline
    2000    & 0,326 & 0,277 & 12,699  & 0,217 \\ \hline
    4000    & 0,472 & 0,315 & 48,507  & 0,447 \\ \hline
    8000    & 0,744 & 0,612 & 189,495 & 0,897 \\ \hline
    16~000  & 1,437 & 1,348 & 779,3   & 2,115 \\ \hline
  \end{tabular}
\end{bbttable}

\begin{bbttable}{Число сравнений символов}{tab:comparisons}
  \begin{tabular}{|r|c|c|c|c|}
    \hline
    & \multicolumn{2}{c|}{Случайный текст} & \multicolumn{2}{c|}{Худший случай} \\ \cline{2-5}
    $n$ & наивный & КМП & наивный & КМП \\ \hline
    1000    & 1302    & 1221    & 47~550     & 1951    \\ \hline
    2000    & 2644    & 2476    & 190~100    & 3901    \\ \hline
    4000    & 5268    & 4893    & 760~200    & 7801    \\ \hline
    8000    & 10~513  & 9872    & 3~040~400  & 15~601  \\ \hline
    16~000  & 21~084  & 19~789  & 12~160~800 & 31~201  \\ \hline
  \end{tabular}
\end{bbttable}

\subsection{Сравнение с оценками сложности (пункт~\ref{task:analysis})}
Зависимость времени от длины текста в логарифмических осях показана на
рисунке~\ref{fig:loglog}. В таких осях степенная зависимость $t \sim n^k$ — прямая с
наклоном $k$: при удвоении $n$ время растёт в $2^k$ раз.

\bbtScreenshot[0.85\linewidth]{materials/results/lab05_loglog.png}{Время поиска от длины
текста, логарифмические оси}{fig:loglog}

На худшем входе время наивного поиска при каждом удвоении $n$ растёт в 3,8–4,5 раза
(таблица~\ref{tab:time}), то есть наклон близок к 2. Это совпадает с оценкой
\eqref{eq:naive}: при $m = n/20$ число сравнений растёт как $n^2$, и в
таблице~\ref{tab:comparisons} оно увеличивается почти ровно в 4 раза на каждом шаге. КМП на
том же входе растёт примерно вдвое на шаг, наклон около 1, а число сравнений ровно
$2n - m + 1$: каждый символ текста сравнивается один раз при совпадении и один раз после
сдвига. При $n = 16~000$ наивный поиск медленнее КМП в 368 раз.

На случайном тексте оба алгоритма растут линейно, а их время почти совпадает: 1,437 и
1,348~мс при $n = 16~000$. Несовпадение со случайным образцом обычно происходит на первом
символе, поэтому наивный алгоритм делает в среднем
$1 + \frac{1}{4} + \frac{1}{16} + \ldots = \frac{4}{3}$ сравнения на позицию, и его худшая
оценка $O(n \cdot m)$ на таких данных не проявляется. Выигрыш КМП виден только на текстах с
длинными частичными совпадениями.

При $n \leq 2000$ время на случайном тексте составляет доли миллисекунды, и точки на
графике отклоняются от прямой сильнее: на таких временах заметны помехи от системы.

\subsection{Проверка корректности (пункт~\ref{task:check})}
Функция \texttt{self\_check} (листинг~\ref{lst:check}) сравнивает результаты обоих алгоритмов
со списком вхождений, найденным методом \texttt{str.find}, на 300 случайных входах над
алфавитом \texttt{ab}, а затем проверяет граничные случаи: пустой текст, образец длиннее
текста и перекрывающиеся вхождения \texttt{aa} в строке \texttt{aaaa}.

\bbtCodeFile[python, lines=99-113]{materials/code/lab05.py}{Самопроверка алгоритмов}{lst:check}

Начало вывода программы приведено в листинге~\ref{lst:output}: расхождений не найдено.

\begin{bbtcode}{text}{Начало вывода программы}{lst:output}
Вариант 7, seed 107, Python 3.12.16
self_check: 300 случайных входов и 3 граничных случая, расхождений с str.find нет
\end{bbtcode}
